%&pdflatex
\documentclass[a4paper,12pt]{article}

\usepackage[spanish]{babel}

\usepackage[leqno]{amsmath}
\usepackage{amsfonts,amsthm}

\nofiles %% .aux files unnecessary

\theoremstyle{plain}
\newtheorem*{propoUno}{Proposici\'on 1}
\newtheorem*{propoDos}{Proposici\'on 2}
\newtheorem*{teo}{Teorema}

\theoremstyle{definition}
\newtheorem*{nota}{Nota}

\newcommand{\parrafo}[1]{\section*{\large #1}}
%% Cambiado tambi\'en el tam\~no de "Referencias" (\bibname)
%% Ver en \begin{thebibliography}

\newcommand{\sop}{\operatorname{sop}}

\newcommand{\N}{\mathbb{N}}
\newcommand{\R}{\mathbb{R}}

\date{}

%%%%%%%%%%%%%%%%%%%%%%%%%%%%%%%%%%%%%%%%%%%%%%%%%%%%%%%%%%%%%%%
\title{\large \bfseries \boldmath CONVERGENCIA EN MEDIA DE LA
SERIE DE FOURIER RESPECTO DE POLINOMIOS ASOCIADOS A LA MEDIDA
$(1-x)^{\alpha} (1+x)^{\beta}\,dx + M \delta_{-1} + N \delta_{1}$
\thanks{Los autores han sido subvencionados por CAICYT PB 85-0338.}}

\author{%
J. J. Guadalupe Hern\'andez$^{(1)}$, M. P\'erez Riera$^{(1)}$,\\[2pt]
F. J. Ruiz Blasco$^{(1)}$ y J. L. Varona Malumbres$^{(2)}$\\[12pt]
\small $^{(1)}$ Departamento de Matem\'aticas,
Universidad de Zaragoza.\\
\small $^{(2)}$ Departamento de Matem\'atica Aplicada,
Universidad de Zaragoza.%
%\renewcommand\thefootnote{}%
\renewcommand\footnotemark{}%
\footnote{Este art\'{\i}culo fue publicado en
``Polinomios Ortogonales y sus Aplicaciones''
(Actas del V Simposium, Vigo, 1988),
91--99, Universidad de Santiago, Vigo, 1989.}%
}
%%%%%%%%%%%%%%%%%%%%%%%%%%%%%%%%%%%%%%%%%%%%%%%%%%%%%%%%%%%%%%%

\begin{document}

%%%%%%%%%%%%%%%%%%%%%%%%%%%%%%%%%%%%%%%%%%%%
\maketitle
%%%%%%%%%%%%%%%%%%%%%%%%%%%%%%%%%%%%%%%%%%%%

%%%%%%%%%%%%%%%%%%%%
\parrafo{0. Introducci\'on}
%%%%%%%%%%%%%%%%%%%%

Sea $d\alpha$ una medida positiva sobre $\R$ para la que podamos
asegurar la existencia de un sistema de polinomios ortonormales
asociado $\{P_{n}\}_{n=0}^{\infty}$. Dada una funci\'on $f$ integrable
para dicha medida, denotaremos por
$$
S_{N}f(x) = \sum_{n=0}^{N} c_{n}P_{n}(x)
= \int_{\R} K_{N}(x,y)f(y)\,d\alpha(y),
$$
donde
$$
K_{N}(x,y) = \sum_{n=0}^{N} P_{n}(x)P_{n}(y)
\quad\text{y}\quad
c_{n} = \int_{\R} f(t)P_{n}(t)\,d\alpha(t),
$$
la suma parcial $N$-\'esima de la serie de Fourier de $f$
respecto de dicho sistema de polinomios.

Siendo claro, por teor\'{\i}a de espacios de Hilbert, que si
$f \in L^{2}(\R;d\alpha)$ dichas sumas parciales convergen
a la funci\'on $f$ en $L^{2}(\R;d\alpha)$, la cuesti\'on natural
que se plantea es: ?`para qu\'e valores de $p$, $1 < p < \infty$,
sigue sucediendo lo mismo? No se conoce ninguna respuesta de
tipo general a esta pregunta y, como consecuencia, las investigaciones
de los matem\'aticos que se han dedicado a este problema han
ido encaminadas a su resoluci\'on para sistemas de polinomios
particulares.

Limit\'andonos a medidas soportadas en el intervalo $[-1,1]$, la
historia del problema se puede resumir en lo siguiente: en 1947,
Pollard~[7] prueba que para polinomios de Legendre, hay convergencia
en $L^{p}([-1,1];dx)$ para $p$ comprendido entre $4/3$ y $4$. El resultado
es extendido por Muckenhoupt~[5] para los polinomios de Jacobi,
probando que para $1 < p < \infty$ hay
convergencia en $L^{p}([-1,1];d\mu)$ si y s\'olo si se cumplen
las ecuaciones
\begin{equation}
\left\{
\begin{aligned}
&4(\alpha+1)^{-1}(2\alpha+1) < p^{-1} < 4(\alpha+1)^{-1}(2\alpha+3),\\
&4(\beta+1)^{-1}(2\beta+1) < p^{-1} < 4(\beta+1)^{-1}(2\beta+3),
\end{aligned}
\right.
\tag{$*$}
\end{equation}
siendo $\alpha,\beta$ los exponentes en $d\mu(x) =
(1-x)^{\alpha}(1+x)^{\beta}\,dx$ ($\alpha,\beta>-1$).
Posteriormente, el problema es completamente resuelto por Badkov
para polinomios de Jacobi generalizados, asociados a ciertas
modificaciones de los pesos de Jacobi por funciones del tipo
$\prod|x-x_{0}|^{\gamma}$. Las t\'ecnicas de demostraci\'on
en~[5] y~[7] pasan por el conocimiento de acotaciones uniformes
para los polinomios junto con propiedades de acotaci\'on de la
transformada de Hilbert y en [1] por una estimaci\'on directa
del tama\~no de los n\'ucleos~$K_{N}$.

El objeto del presente art\'{\i}culo es resolver el problema anterior
para medidas que son modificaciones por deltas de Dirac en los
extremos de los pesos de Jacobi, esto es,
$d\nu(x) = (1-x)^{\alpha}(1+x)^{\beta}\,dx
+ M\delta_{-1} + N\delta_{1}$ en $[-1,1]$. Algunos aspectos
de los polinomios asociados a dichas medidas han sido estudiados
en~[2] y por este motivo nos referiremos a ellos como polinomios
de Koornwinder.

Demostraremos que el hecho de a\~nadir deltas de Dirac en los
extremos no modifica el rango de los $p$ para los que hay convergencia,
es decir, habr\'a convergencia cuando y s\'olo cuando $p$ satisfaga
las ecuaciones~($*$).

Utilizaremos las siguientes notaciones:
$$
d\mu(x) = (1-x)^{\alpha}(1+x)^{\beta}\,dx
\quad\text{sobre}\quad [-1,1] \qquad(\alpha,\beta > -1);
$$
$$
d\nu(x) = (1-x)^{\alpha}(1+x)^{\beta}\,dx + M\delta_{-1} + N\delta_{1}
\quad\text{sobre}\quad [-1,1].
$$
$P_{n}(x) = k_{n}x^{n} + \cdots$, ser\'an, en general,
los polinomios ortonormales asociados a una medida $d\alpha$
y escribiremos $P_{n}^{a} = k_{n}^{a} x^{n} + \cdots$,
para denotar los asociados a la medida $(x-a)^{2}\,d\alpha(x)$.

El art\'{\i}culo est\'a organizado como sigue: en el primer p\'arrafo
obtenemos una f\'ormula de tipo general que relaciona los polinomios
asociados a una modificaci\'on de un peso por una delta de Dirac
en funci\'on de los de partida y de los asociados a la modificaci\'on
por $(x-a)^{2}$. Esta f\'ormula nos permitir\'a en el segundo apartado
obtener acotaciones para los polinomios de Koornwinder y, por
\'ultimo, en el p\'arrafo~3, demostramos nuestro resultado sobre
convergencia.

%%%%%%%%%%%%%%%%%%%%
\parrafo{1. Polinomios asociados a una modificaci\'on por una
delta de Dirac}
%%%%%%%%%%%%%%%%%%%%

\begin{propoUno}
Sean $d\alpha$ una medida positiva sobre $\R$ y $a \in \R$,
y denotemos por $Q_{n}$ la sucesi\'on de polinomios
ortonormales con respecto a $d\alpha + M\delta_{a}$,
siendo $M > 0$ y $\delta_{a}$ la delta de Dirac en~$a$. Entonces,
para todo $n$ existen $A_{n}, B_{n} \in (0,1)$ tales que
$$
Q_{n}(x) = A_{n}P_{n}(x) + B_{n}(x-a)P_{n-1}^{a}(x).
$$
Adem\'as, si $\sop d\alpha = [-1,1]$ , $\alpha' > 0$
a.e. y $a \in [-1,1]$, se tiene:
$$
\lim_{n\to\infty} A_{n}K_{n-1}(a,a) = [\lambda(a) + M]^{-1},
\quad \lim_{n\to\infty} B_{n} = M [\lambda(a) + M]^{-1},
$$
donde
$$
\lambda(a) = \lim_{n\to\infty} K_{n}(a,a)^{-1} \in [0,\infty).
$$
\end{propoUno}

\begin{proof}
Imponiendo que $P_{n}(x) + C_{n}(x-a)P_{n-1}^{a}(x)$
sea ortogonal al espacio de los polinomios de grado menor o igual
que $n-1$, con respecto a $d\alpha + M\delta_{a}$ (utilizando
la base $\{(x-a)^{j}\}$) se obtiene que la constante $C_{n}$ debe ser
$M(k_{n}/k_{n-1}^{a})K_{n-1}(a,a)$. As\'{\i}, dividiendo dicho polinomio
por su norma, obtenemos $Q_{n}$, dado que $P_{n}(x)+C_{n}(x-a)P_{n-1}^{a}(x)$
es un polinomio de grado exactamente $n$ y con coeficiente director
positivo. Su norma en $L^{2}(d\alpha + M\delta_{a})$ es
$$
D_{n} = [ M P_{n}(a)^{2} + 1 + Cs(2,n) + 2 C_{n} (k_{n}/k_{n-1}^{a}) ]^{1/2}.
$$
Por tanto, $Q_{n}(x) = D_{n}^{-1} P_{n}(x) + C_{n} D_{n}^{-1} (x-a)P_{n-1}^{a}(x)$.
Como $D_{n} > \max\{1,C_{n}\}$,
se sigue que $A_{n} = D_{n}^{-1}, B_{n} = C_{n} D_{n}^{-1} \in (0,1)$.

Para la segunda parte de la proposici\'on basta tener en cuenta
la expresi\'on de $D_{n}$ y $C_{n}$ y recordar los siguientes resultados
de tipo general:

--- Siempre que $\sop d\alpha = [-1,1]$ y $\alpha' > 0$
a.e., se tiene que $(k_{n}/k_{n+1}) \to 1/2$
(ver~[9] , p\'ag.~212).

--- En las mismas hip\'otesis,
$$
(k_{n}^{a}/k_{n}) \to \exp \left(\frac{-1}{4\pi}
\int_{0}^{2\pi} \log\{(\cos t-a)^{2}\}\,dt\right)
$$
(ver~[4], teorema~11) y
puede calcularse que, $\forall a \in [-1,1]$, el valor de esta
integral es $-4\pi \log 2$. Por lo tanto,
$(k_{n}/k_{n-1}^{a}) \to 1$.

--- Por otro lado, de $\sop d\alpha = [-1,1]$ y $\alpha' > 0$ a.e.,
se sigue que, en la relaci\'on de recurrencia $x P_{n}
= a_{n+1}P_{n+1} + b_{n} P_{n} + a_{n} P_{n-1}$, $a_{n} \to 1/2$ y $b_{n} \to 0$
(ver~[9], p\'ag.~212, o bien~[4], teorema~10). Y entonces se
tiene:
$$
\lim_{n\to\infty} [P_{n}(x)^{2}/K_{n-1}(x,x)] = 0
\quad\forall x \in [-1,1]
$$
(ver~[6], teorema~3, p\'ag.~26).
Por consiguiente, $\lim_{n\to\infty} [P_{n}(a)/K_{n-1}(a,a)] = 0$.
\end{proof}

%%%%%%%%%%%%%%%%%%%%
\parrafo{2. Estimaciones para los polinomios de Koornwinder}
%%%%%%%%%%%%%%%%%%%%

\begin{propoDos}
Sea $\{Q_{n}\}$ la sucesi\'on de polinomios
ortonormales con respecto a la medida $d\nu(x) = (1-x)^{\alpha}(1+x)^{\beta}\,dx
+ M\delta_{-1}(x) + N\delta_{1}(x$) sobre $[-1,1]$, $M,N \geq 0$.
Entonces existe una constante $C$ tal que $\forall x\in [-1,1]$
y $\forall n\in \N$:
\begin{itemize}
\item[(a)] $|Q_{n}(x)| \leq C (1-x + n^{-2})^{-\alpha/2-1/4}
(1+x + n^{-2})^{-\beta/2-1/4}$.

\item[(b)] Si $M>0$, $|Q_{n}(-1)| \leq C n^{-\beta-3/2}$;
si $N>0$, $|Q_{n}(1)| \leq C n^{-\alpha-3/2}$.
\end{itemize}
\end{propoDos}

\begin{proof}
(a) Sean $P_{n}$ los polinomios ortonormales con respecto a
$d\mu = (1-x)^{\alpha}(1+x)^{\beta}\,dx$,
$R_{n}$ los ortonormales con respecto a $d\mu + M\delta_{-1}$
y $P_{n}^{(-1)}$ los ortonormales con respecto a $(1+x)^{2}\,d\mu$.
Puesto que $d\mu$ y $(1+x)^{2}\,d\mu$ son medidas de Jacobi,
se verifican las acotaciones
$$
|P_{n}(x)| \leq C (1-x + n^{-2})^{-\alpha/2-1/4}
(1+x + n^{-2})^{-\beta/2-1/4},
$$
$$
|P_{n-1}^{(-1)}(x)| \leq C (1-x + n^{-2})^{-\alpha/2-1/4}
(1+x + n^{-2})^{-\beta/2-5/4}
$$
(ver [5]). De la proposici\'on~1 se deduce ahora que
$$
|R_{n}(x)| \leq C (1-x + n^{-2})^{-\alpha/2-1/4}
(1+x + n^{-2})^{-\beta/2-1/4}.
$$
Repitiendo el proceso con la medida $d\mu+M\delta_{-1}$
y $a = 1$ se llega a la acotaci\'on de los~$Q_{n}$.

(b) Si $\{P_{n}\}$ y $\{K_{n}\}$ son las sucesiones de polinomios ortonormales
y n\'ucleos de $d\mu$, es conocido (ver~[10], {\S}\,4.1, {\S}\,4.3
y {\S}\,4.5) que $|P_{n}(-1)| \sim n^{\beta+1/2}$
y $|K_{n}(-1,-1)| \sim n^{2(\beta+1)}$,
donde con $a_{n}\sim b_{n}$ indicamos que
existen $C,D > 0$ tales que, $\forall n$,
$C b_{n} \leq a_{n} \leq D b_{n}$.
De esto se deduce por la proposici\'on~1 que, si $M > 0$,
los polinomios $R_{n}$ ortonormales con respecto a $d\mu+M\delta_{-1}$
cumplen
$$
|R_{n}(-1)| = A_{n}|P_{n}(-1)| \sim |P_{n}(-1)/K_{n}(-1,-1)|
\leq Cn^{-\beta-3/2}.
$$
Por la misma raz\'on, los polinomios $R_{n}^{1}$, ortonormales
con respecto a $(1-x)^{2}(d\mu+M\delta_{-1}) = (1-x)^{2}\,d\mu+4M\delta_{-1}$
cumplen que $|R_{n}^{1}(-1)| \leq Cn^{-\beta-3/2}$.
De nuevo por la proposici\'on~1,
$Q_{n}(-1) = A_{n}R_{n}(-1) - 2 B_{n}R_{n-1}^{1})(-1)$
y por tanto $|Q_{n}(-1)| \leq Cn^{-\beta-3/2}$.
An\'alogamente, $|Q_{n}(1)| \leq C n^{-\alpha-3/2}$
si $N > 0$.
\end{proof}

%%%%%%%%%%%%%%%%%%%%
\parrafo{3. Convergencia de la serie de Fourier}
%%%%%%%%%%%%%%%%%%%%

\begin{teo}
Denotando por $S_{n}f$ las sumas parciales del desarrollo
en serie de Fourier de $f$ con respecto a $d\nu(x) = (1-x)^{\alpha}(1+x)^{\beta}\,dx
+ M\delta_{-1} + N\delta_{1}$ sobre $[-1,1]$, las siguientes
afirmaciones son equivalentes:
\begin{itemize}
\item[(i)] $S_{n}$ son operadores uniformemente acotados en $L^{p}(d\nu)$.

\item[(ii)] $S_{n}f \to f$ en $L^{p}(d\nu)$, para toda funci\'on $f$ en $L^{p}(d\nu)$.

\item[(iii)] $p$ satisface las condiciones~($*$).
\end{itemize}
\end{teo}

\begin{proof}
La equivalencia entre (i) y (ii) es cl\'asica para cualquier
medida finita sobre $[-1,1]$, sigui\'endose por bien conocidos
argumentos de An\'alisis Funcional (teorema de la acotaci\'on
uniforme y densidad de los polinomios en $L^{p}(d\nu)$). Por
otra parte, (i)~$\Rightarrow$~(iii) consiste en la simple observaci\'on
de que el teorema de M\'at\'e, Nevai y Totik referente a condiciones
necesarias para la convergencia de la serie de Fourier (ver~[3])
solamente requiere el conocimiento de la parte absolutamente
continua de la medida y, por tanto, el rango de los $p$ es el mismo
que para los correspondientes polinomios de Jacobi. Por tanto,
s\'olo debemos ocuparnos de la implicaci\'on
(iii)~$\Rightarrow$~(i).

Se trata de probar que la expresi\'on
$$
\int_{-1}^1 |S_{n}f(x)|^{p}\,d\nu(x) = M|S_{n}f(-1)|^{p}
+ N|S_{n}f(1)|^{p} + \int_{-1}^1 |S_{n}f(x)|^{p}\,d\mu(x),
$$
est\'a acotada (independientemente de $n$) por
$C \int_{-1}^1 |f(x)|^{p}\,d\nu(x)$.
Adem\'as, basta hacerlo para $p > 2$,
por dualidad.

Para la acotaci\'on del primer sumando, observemos que
$$
|S_{n}f(-1)| = \left|\sum_{j=0}^n \widehat{f}(j)P_{j}(-1)\right|
\leq \left\{\sum_{j=0}^n \widehat{f}(j)^{2}\right\}^{1/2}
\left\{\sum_{j=0}^n P_{j}(-1)^{2}\right\}^{1/2}.
$$
El primer factor es menor o igual que $\|f\|_{L^2(d\nu)}
\leq C \|f\|_{L^p(d\nu)}$ y el
segundo est\'a acotado uniformemente por una constante, debido
a la parte~(b) de la proposici\'on~2.
La cota para $|S_{n}f(1)|$ es totalmente an\'aloga.

Para la acotaci\'on del tercer sumando, tenemos que
\begin{multline*}
\int_{-1}^1 |S_{n}f(x)|^{p}\,d\mu(x) \\
\leq C M^{p}|f(-1)|^{p} \int_{-1}^1 |K_{n}(x,-1)|^{p}\,d\mu(x)
+ C N^{p}|f(1)|^{p} \int_{-1}^1 |K_{n}(x,1)|^{p}\,d\mu(x) \\
\null+ C \int_{-1}^1 \left| \int_{-1}^1 K_{n}(x,y)f(y)\,d\mu(y) \right|^{p}
\,d\mu(x).
\end{multline*}
\sloppy Como $M |f(-1)|^{p}$, $N |f(1)|^{p}$ y
$\int_{-1}^1 |f(x)|^{p}\,d\mu(x)$ est\'an controlados
por $\int_{-1}^1 |f(x)|^{p}\,d\nu(x)$, para llegar
a la demostraci\'on de (i) basta que logremos las tres siguientes
mayoraciones:
\begin{align}
\int_{-1}^1 |K_{n}(x,-1)|^{p}\,d\mu(x) &\leq C,
\quad\text{para todo } n,
\tag{A-1}\\
\int_{-1}^1 |K_{n}(x,1)|^{p}\,d\mu(x) &\leq C,
\quad\text{para todo } n,
\tag{A-2}\\
\int_{-1}^1 \left| \int_{-1}^1 K_{n}(x,y)f(y)\,d\mu(y) \right|^{p}\,d\mu(x)
&\leq C \int_{-1}^1 |f(x)|^{p}\,d\mu(x),
\quad\text{para todo } n.
\tag{A-3}
\end{align}
Para la demostraci\'on de (A-1) (la de (A-2) es similar), partimos
de la descomposici\'on del n\'ucleo (ver [8]) como
\begin{multline*}
K_{n}(x,y) = r_{n}Q_{n+1}(x)Q_{n+1}(y) \\
+ s_{n} \frac{(1-x^{2})R_{n}(x)Q_{n+1}(y)}{x-y}
- s_{n} \frac{(1-y^{2})R_{n}(y)Q_{n+1}(x)}{x-y}\,,
\end{multline*}
donde $\{r_{n}\}$ y $\{s_{n}\}$ son sucesiones acotadas (ver~[11]),
$Q_{n}(x)$ son los polinomios ortonormales con respecto a $d\nu(x)$
y $R_{n}(x)$ son los polinomios ortonormales correspondientes a
la medida $(1-x^{2})\,d\nu(x) = (1-x^{2})\,d\mu(x)$. Entonces,
para $y = 1$, resulta
$K_{n}(x,1) = r_{n}Q_{n+1}(x)Q_{n+1}(1) - s_{n}(1+x)R_{n}(x)Q_{n+1}(1)$.
A partir de esta expresi\'on, con las cotas obtenidas para $Q_{n}$
en la proposici\'on~2 y las acotaciones que verifican los $R_{n}$
(puesto que son polinomios de Jacobi con exponentes $\alpha+1$
y $\beta+1$), es un sencillo ejercicio obtener la aco\-ta\-ci\'on~(A-2).

\fussy
Por \'ultimo, utilizando la misma descomposici\'on del n\'ucleo,
la obtenci\'on de (A-3) consiste en demostrar que
$$
\int_{-1}^1 |W_{i,n}f(x)|^{p}\,d\mu(x) \leq C
\int_{-1}^1 |f(x)|^{p}\,d\mu(x) , \quad\forall n,
\quad\text{para } i = 1,2,3,
$$
donde
$$
W_{i,n}f(x) = \int_{-1}^1 T_{i}(n,x,y)f(y)\,d\mu(y)
$$
y
\begin{align*}
T_{1}(n,x,y) &= Q_{n+1}(x)Q_{n+1}(y) , \\
T_{2}(n,x,y) &= \frac{(1-x^{2})R_{n}(x)Q_{n+1}(y)}{x-y}\,, \\
T_{3}(n,x,y) &= \frac{(1-y^{2})R_{n}(y)Q_{n+1}(x)}{x-y}\,.
\end{align*}
Pero esta demostraci\'on es id\'entica a la del caso de los pesos
de Jacobi, sin deltas de Dirac, puesto que en \'esta s\'olo se
necesitan cotas superiores para los polinomios ortonormales,
que, seg\'un la proposici\'on~2, son v\'alidas tambi\'en para
los polinomios de Koornwinder. Esto concluye la demostraci\'on
de (iii)~$\Rightarrow$~(i) y del teorema.
\end{proof}

\begin{nota}
Siguiendo en parte el esquema anterior, puede estudiarse
el problema an\'alogo para los polinomios ortogonales con respecto
a medidas sobre el intervalo $[-1,1]$ que son suma de un peso de
Jacobi generalizado y de deltas de Dirac incluso en puntos del
interior del intervalo. Este problema ser\'a objeto de un posterior
trabajo.
\end{nota}



%%%%%%%%%%%%%%%%%%%%%%%%%%%%%%%%%%%
\renewcommand{\refname}{\large Referencias}
\begin{thebibliography}{11}
%%%%%%%%%%%%%%%%%%%%%%%%%%%%%%%%%%%

\frenchspacing

\bibitem{1.}
\textsc{V. M. Badkov},
Convergence in the mean and almost everywhere
of Fourier series in polynomials orthogonal on an interval,
\textit{Math. USSR Sbornik} \textbf{24} (1974), 223--256.

\bibitem{2.}
\textsc{T. H. Koornwinder},
Orthogonal polynomials with weight function
$(1-x)^{\alpha} (1+x)^{\beta} + M \delta(x+1) + N \delta(x-1)$,
\textit{Canad. Math. Bull.} \textbf{27} (1984), 205--214.

\bibitem{3.}
\textsc{A. M\'at\'e, P. Nevai y V. Totik},
Necessary conditions for weighted mean convergence of Fourier
series in orthogonal polynomials,
\textit{J. Approx. Theory} \textbf{46} (1986), 314--322.

\bibitem{4.}
\textsc{A. M\'at\'e, P. Nevai y V. Totik},
Extensions of Szeg\H{o}'s theory of orthogonal polynomials.~II,
\textit{Constr. Approx.} \textbf{3} (1987), 51--72.

\bibitem{5.}
\textsc{B. Muckenhoupt},
Mean convergence of Jacobi series,
\textit{Proc. Amer. Math. Soc.} \textbf{23} (1969), 306--310.

\bibitem{6.}
\textsc{P. Nevai}, ``Orthogonal Polynomials'',
Memoirs of the Amer. Math. Soc. 213 (1979).

\bibitem{7.}
\textsc{H. Pollard},
The mean convergence of orthogonal series~I,
\textit{Trans. Amer. Math. Soc.} \textbf{62} (1947), 387--403.

\bibitem{8.}
\textsc{H. Pollard},
The mean convergence of orthogonal series~II,
\textit{Trans. Amer. Math. Soc.} \textbf{63} (1948), 355--367.

\bibitem{9.}
\textsc{E. A. Rahmanov},
On the asymptotics of the ratio of orthogonal polynomials,
\textit{Math. USSR Sbornik} \textbf{32} (1977), 199--213.

\bibitem{10.}
\textsc{G. Szeg\H{o}},
``Orthogonal Polynomials'',
Amer. Math. Soc. Colloq. Publ., vol.~23,
Amer. Math. Soc., Providence, R.I., 1959.

\bibitem{11.}
\textsc{J. L. Varona},
``Convergencia de series de Fourier respecto de sistemas ortogonales'',
Publicaciones del Seminario Matem\'atico Garc\'{\i}a
de Galdeano, Zaragoza, 1987.

\end{thebibliography}

\end{document}
